\documentclass[../main.tex]{subfiles}
\begin{document}
\chapter{2015-2016学年微积分（一）（上）期中考试}

\section{基本计算题(每小题 5 分，共 60 分)}
\begin{enumerate}
    \item 计算数列极限$\lim_{n\to\infty}\left(\frac{1}{n^{2}+n+1}+\frac{2}{n^{2}+n+2}
              +\cdots+\frac{n}{n^{2}+n+n}\right)$.
          \vspace{7em}

    \item 计算极限$l=\lim_{x\to 0}\left(\frac{\mathrm{e}^{x}+\mathrm{e}^{2x}+\cdots+\mathrm{e}^{nx}}{n}
              \right)^{\frac{1}{x}}$.
          \vspace{7em}

    \item 计算极限$l=\lim_{x\to +\infty}x^{2}\ln\left(x\sin\frac{1}{x}\right)$.
          \vspace{7em}

    \item 已知$\lim_{x\to+\infty}\left(\frac{x^{2}+x+1}{1-x}-ax-b\right)=0$，求常数$a
              ,b$的值.
          \vspace{8em}

    \item 求当$x\to0$时，无穷小$u(x)=\arcsin x-x$的主部与阶数.
          \vspace{8em}

    \item 设函数$y=\ln\sqrt{\frac{2}{\pi}\arctan\frac{1}{x}}$，求$y'(1)$.
          \vspace{8em}

    \item 设函数$y=\frac{\sin x}{x}$，求$\frac{\mathrm{d}y}{\mathrm{d}x},\frac{\mathrm{d}y}{\mathrm{d}(\cot
                  x)}$.
          \vspace{7em}

    \item 设函数$f(x)=\frac{1}{2^{x}}-x^{3}$，$x=\varphi(y)$是$y=f(x)$的反函数，求$\left
              .\varphi'(y)\right|_{y=1}$.
          \vspace{7em}

    \item 设函数$y=\frac{1}{2x^{2}-3x+1}$，求$y^{(10)}(x)$.
          \vspace{8em}

    \item 设函数$y=y(x)$由参数方程$\begin{cases}
                  x=\ln(1+t^{2}), \\
                  y=t-\arctan t
              \end{cases}$所确定，求$\frac{\mathrm{d}^{2}y}{\mathrm{d}x^{2}}$.
          \vspace{8em}

    \item 求函数$f(x)=\frac{x\sin(1-x)}{|x|(x^{2}-1)}$的间断点，并判断其类型.
          \vspace{7em}

    \item 一架巡逻直升机在距地面$3\mathrm{km}$的高度以$120\mathrm{km/h}$匀速地沿一条水平笔直的高速公路向前飞行，飞行员观察到迎面驶来一辆汽车.设汽车行进的速度为匀速，当直升机与汽车间的距离为$5
              \mathrm{km}$时通过雷达测出此距离以$160\mathrm{km/h}$的速率减小，试求汽车行进的速度.
          \vspace{7em}
\end{enumerate}

\section{综合题(每小题 6 分，共 30 分)}
\begin{enumerate}
    \setcounter{enumi}{12}
    \item 设函数$f(x)=
              \begin{cases}
                  x^{2}\sin\frac{1}{x}, & x\neq0, \\
                  0,                    & x=0
              \end{cases}$，求$f'(x)$，并讨论$f'(x)$的连续性.
          \vspace{8em}

    \item 设$y=y(x)$由方程$y=\sin(x+y)$确定，求$\frac{\mathrm{d}^{2} y}{\mathrm{d}x^{2}}$.
          \vspace{9em}

    \item 设函数$f(x)$满足$f'(0)=0$，$f''(0)$存在，求极限$\lim_{x\to0}\frac{f(x)-f(\ln(1+x))}{x^{3}}$.
          \vspace{9em}

    \item 设$f(x)$是周期为5的连续函数，在$x=1$处可导，在$x=0$附近满足
          \[
              f(1+\sin x)-3f(1-\sin x)=8x+o(x),
          \]
          求曲线$y=f(x)$在点$(6, f(6))$处的切线方程.
          \vspace{10em}

    \item 设$x_{1}=2$，$x_{n}=2+\frac{1}{x_{n-1}}(n>1)$，证明$\lim_{n\to\infty}
              x_{n}$存在，并求其值.
          \vspace{10em}
\end{enumerate}

\section{证明题(每小题 5 分，共 10 分)}
\begin{enumerate}
    \setcounter{enumi}{17}
    \item 设$f(x)$在闭区间$[a,b]$上有二阶导数，且$f(a)=f(b)=0$，$f'_{+}(a)\cdot
              f'_{-}(b)>0$，证明：

          （1）至少存在一点$\xi\in(a,b)$使得$f(\xi)=0$；

          （2）至少存在一点$\eta\in(a,b)$使得$f''(\eta)=0$.
          \vspace{10em}

    \item 设函数$f(x)$在区间$[0,1]$上具有二阶导数，且有正常数$a,b$，使得$|f(
              x)|\leq a$，$|f''(x)|\leq b$.证明：对任意$c\in(0,1)$，有$|f'(c)|\leq 2a+\frac{b}{2}$.
          \vspace{10em}
\end{enumerate}

\chapter{2015-2016学年微积分（一）（上）期中考试参考答案}
\section{基本计算题(每小题 5 分，共 60 分)}
\begin{enumerate}
    \item \textbf{Solution}. 设$x_{n}=\frac{1}{n^{2}+n+1}+\frac{2}{n^{2}+n+2}+\cdots
              +\frac{n}{n^{2}+n+n}$，则
          \[
              \frac{1+2+\cdots+n}{n^{2}+n+n}< x_{n} < \frac{1+2+\cdots+n}{n^{2}+n+1},
          \]

          又$\lim_{n\to\infty}\frac{1+2+\cdots+n}{n^{2}+n+n}=\lim_{n\to\infty}\frac{n(n+1)}{2(n^{2}+2n)}
              =\frac{1}{2}$,

          $\lim_{n\to\infty}\frac{1+2+\cdots+n}{n^{2}+n+1}=\lim_{n\to\infty}\frac{n(n+1)}{2(n^{2}+n+1)}
              =\frac{1}{2}$,

          由夹逼定理得$\lim_{n\to\infty}x_{n}=\frac{1}{2}$.

    \item \textbf{Solution}.
          \[
              \begin{aligned}
                  l = & \exp\left[\lim_{x\to0}\frac{1}{x}\ln\left(\frac{\mathrm{e}^x+\mathrm{e}^{2x}+\cdots+\mathrm{e}^{nx}}{n}\right)\right]                                           \\
                  =   & \exp\left[\lim_{x\to0}\frac{1}{x}\left(\frac{\mathrm{e}^x+\mathrm{e}^{2x}+\cdots+\mathrm{e}^{nx}}{n}-1\right)\right]                                            \\
                  =   & \exp\left[\lim_{x\to0}\frac{1}{n}\left(\frac{\left(\mathrm{e}^x-1\right)+\left(\mathrm{e}^{2x}-1\right)+\cdots+\left(\mathrm{e}^{nx}-1\right)}{x}\right)\right] \\
                  =   & \exp\left[\frac{1}{n}(1+2+\cdots+n)\right]=\mathrm{e}^{\frac{n+1}{2}}.
              \end{aligned}
          \]

    \item \textbf{Solution}. 令$x=\frac{1}{t}$，则
          \[
              \begin{aligned}
                  l= & \lim_{x\to+\infty}\frac{\ln\left(x\sin\frac{1}{x}\right)}{\frac{1}{x^2}}= \lim_{t\to0^+}\frac{\ln\left(\frac{\sin t}{t}\right)}{t^2} \\
                  =  & \lim_{t\to0^+}\frac{\frac{\sin t}{t}-1}{t^2}=\lim_{t\to0^+}\frac{\sin t-t}{t^3}                                                      \\
                  =  & \lim_{t\to0^+}\frac{\cos t-1}{3t^2}=-\frac{1}{6}.
              \end{aligned}
          \]

    \item \textbf{Solution}.原极限式可变形为
          \[
              0=\lim_{x\to+\infty}\frac{(1+a)x^{2}+(1-a+b)x+1-b}{1-x},
          \]

          得$1+a=0,1-a+b=0$，即$a=-1,b=-2$.

    \item \textbf{Solution}. 设$\lim_{x\to0}\frac{u(x)}{cx^{r}}=1$，则
          \[
              \begin{aligned}
                  1= & \lim_{x\to0}\frac{\arcsin x- x}{cx^{r}}=\lim_{x\to0}\frac{\frac{1}{\sqrt{1-x^{2}}}-1}{crx^{r-1}} \\
                  =  & \lim_{x\to0}\frac{1}{\sqrt{1-x^2}}\cdot\frac{1-\sqrt{1-x^2}}{crx^{r-1}}                          \\
                  =  & \lim_{x\to0}\frac{\frac{1}{2}x^2}{crx^{r-1}},
              \end{aligned}
          \]

          因此必须$\begin{cases}
                  r-1=2, \\
                  cr=\frac{1}{2},
              \end{cases}$，解得$r=3,c=\frac{1}{6}$.

          所以$u(x)$的主部为$\frac{1}{6}x^{3}$，阶数为3.

    \item \textbf{Solution}. $y=\frac{1}{2}\left[\ln\frac{2}{\pi}+\ln\arctan\frac{1}{x}
                  \right]$，

          所以
          \[
              \begin{aligned}
                  y' & =\frac{1}{2}\cdot\frac{1}{\arctan\frac{1}{x}}\cdot\frac{1}{1+\frac{1}{x^{2}}}\cdot\left(-\frac{1}{x^{2}}\right) \\
                     & = -\frac{1}{2(x^{2}+1)\arctan\frac{1}{x}}.
              \end{aligned}
          \]

          代入$x=1$得$y'(1)=-\frac{1}{4\arctan 1}=-\frac{1}{\pi}$.

    \item \textbf{Solution}. $\frac{\mathrm{d}y}{\mathrm{d}x}=\frac{x\cos x-\sin x}{x^{2}}$，

          $\frac{\mathrm{d}y}{\mathrm{d}(\cot x)}=\frac{\frac{\mathrm{d}y}{\mathrm{d}x}}{\frac{\mathrm{d}(\cot
                      x)}{\mathrm{d}x}}=\frac{x\cos x-\sin x}{-x^{2}\csc^{2} x}=\frac{\sin^{3} x-x\sin^{2}
                  x\cos x}{x^{2}}$.

    \item \textbf{Solution}. 当$x=0$时，$f(0)=\frac{1}{2^{0}}-0^{3}=1$.所以
          \[
              \begin{aligned}
                  \left.\varphi'(y)\right|_{y=1} & =\frac{1}{f'(0)}                                                \\
                                                 & = \frac{1}{\left.\frac{1}{2^x}\ln\frac{1}{2}+3x^2\right|_{x=0}} \\
                                                 & = \frac{1}{\ln\frac{1}{2}}=-\frac{1}{\ln 2}.
              \end{aligned}
          \]

    \item \textbf{Solution}. $y=\frac{1}{x-1}-\frac{2}{2x-1}$，所以
          \[
              \begin{aligned}
                  y^{(10)}(x) = & \frac{(-1)^{10}\,10!}{(x-1)^{11}}-\frac{2\cdot(-1)^{10}\,10!\cdot 2^{10}}{(2x-1)^{11}} \\
                  =             & \frac{10!}{(x-1)^{11}}-\frac{2^{11}\cdot 10!}{(2x-1)^{11}}.
              \end{aligned}
          \]

    \item \textbf{Solution}. $\frac{\mathrm{d}y}{\mathrm{d}x}=\frac{\frac{\mathrm{d}y}{\mathrm{d}t}}{\frac{\mathrm{d}x}{\mathrm{d}t}}
              =\frac{1-\frac{1}{1+t^{2}}}{\frac{2t}{1+t^{2}}}=\frac{t}{2}$，

          $\frac{\mathrm{d}^{2}y}{\mathrm{d}x^{2}}=\frac{\mathrm{d}}{\mathrm{d}t}\left
              (\frac{\mathrm{d}y}{\mathrm{d}x}\right)\cdot\frac{\mathrm{d}t}{\mathrm{d}x}
              =\left(\frac{t}{2}\right)'\cdot\frac{1+t^{2}}{2t}=\frac{1+t^{2}}{4t}$.

    \item \textbf{Solution}. 函数$f(x)$的间断点为$x=0,\pm1$.

          当$x=0$时，$\lim_{x\to 0^+}f(x)=\lim_{x\to0^+}\frac{x\sin(1-x)}{x(x^{2}-1)}
              =-\sin 1$，$\lim_{x\to 0^-}f(x)=\lim_{x\to0^-}\frac{x\sin(1-x)}{-x(x^{2}-1)}
              =\sin 1$，

          所以$x=0$为跳跃间断点.

          当$x=1$时，$\lim_{x\to1}f(x)=\lim_{x\to1}\frac{x\sin(1-x)}{x(x-1)(x+1)}=\lim
              _{x\to1}\frac{1-x}{(x-1)(x+1)}=-\frac{1}{2}$，

          且函数在$x=1$处无定义，所以$x=1$为可去间断点.

          当$x=-1$时，$\lim_{x\to-1}f(x)=\lim_{x\to-1}\frac{x\sin(1-x)}{-x(x-1)(x+1)}
              =\infty$，所以$x=-1$为无穷间断点.

    \item \textbf{Solution}. 设$x(t)$为$t$时刻飞机与汽车的水平距离，$y(t)$为$t$时刻飞机与汽车的距离，则
          \[
              x^{2}(t)+h^{2}=y^{2}(t),
          \]

          其中$h=3\mathrm{km}$，方程两边对$t$求导得
          \[
              x\frac{\mathrm{d}x}{\mathrm{d}t}=y\frac{\mathrm{d}y}{\mathrm{d}t},
          \]

          由题设知，在$t=t_{0}$时刻，$y(t_{0})=5\mathrm{km}$，$\left.\frac{\mathrm{d}y}{\mathrm{d}t}
              \right|_{t=t_0}=-160\mathrm{km/h}$，$x(t_{0})=4\mathrm{km}$，

          代入上式得$\left.\frac{\mathrm{d}x}{\mathrm{d}t}\right|_{t=t_0}=\frac{y(t_{0})}{x(t_{0})}
              \cdot\left.\frac{\mathrm{d}y}{\mathrm{d}t}\right|_{t=t_0}=-200\mathrm{km/h}$，

          所以汽车行进的速度为$80\mathrm{km/h}$.
\end{enumerate}

\section{综合题(每小题 6 分，共 30 分)}
\begin{enumerate}
    \setcounter{enumi}{12}
    \item \textbf{Solution}.当$x\neq0$时，$f'(x)=2x\sin\frac{1}{x}-\cos\frac{1}{x}$.

          当$x=0$时，$f'(0)=\lim_{x\to0}\frac{f(x)-f(0)}{x-0}=\lim_{x\to0}\frac{x^{2}\sin\frac{1}{x}}{x}
              =0$，

          所以$f'(x)=
              \begin{cases}
                  2x\sin\frac{1}{x}-\cos\frac{1}{x}, & x\neq0, \\
                  0,                                 & x=0.
              \end{cases}$

          当$x\neq0$时，$f'(x)$显然连续.又$\lim_{x\to0}f'(x)=\lim_{x\to0}\left(2x\sin
              \frac{1}{x}-\cos\frac{1}{x}\right)$不存在，

          所以$f'(x)$在$x=0$处不连续.综上所述，$f'(x)$在$\mathbf{R}\setminus\{0\}$处连续.

    \item \textbf{Solution}. 在方程$y=\sin(x+y)$两边对$x$求导得
          \[
              \frac{\mathrm{d}y}{\mathrm{d}x}=\cos(x+y)\cdot\left(1+\frac{\mathrm{d}y}{\mathrm{d}x}
              \right),
          \]

          解得$\frac{\mathrm{d}y}{\mathrm{d}x}=\frac{\cos(x+y)}{1-\cos(x+y)}$.

          在方程$y'=\cos(x+y)(1+y')$两边再对$x$求导得
          \[
              \frac{\mathrm{d}^{2}y}{\mathrm{d}x^{2}}= -\sin(x+y)\cdot\left(1+\frac{\mathrm{d}y}{\mathrm{d}x}
              \right)^{2}+\cos(x+y)\cdot\frac{\mathrm{d}^{2}y}{\mathrm{d}x^{2}}
          \]

          解得$\frac{\mathrm{d}^{2}y}{\mathrm{d}x^{2}}=\frac{-\sin(x+y)(1+y')^{2}}{1-\cos(x+y)}
              =\frac{-\sin(x+y)}{[1-\cos(x+y)]^{3}}$.

    \item \textbf{Solution}. 由于$f''(0)$存在，则$f(x)$在$x=0$的邻域内可导.

          由Lagrange中值定理，存在$\xi$介于$x$与$\ln(1+x)$之间，使得
          \[
              \begin{aligned}
                  \lim_{x\to0}\frac{f(x)-f(\ln(1+x))}{x^3} & =\lim_{x\to0}\frac{f'(\xi)(x-\ln(1+x))}{x^3}                                           \\
                                                           & =\lim_{x\to0}\frac{f'(\xi)-f'(0)}{\xi-0}\cdot\frac{\xi}{x}\cdot\frac{x-\ln(1+x)}{x^2}.
              \end{aligned}
          \]

          当$x\to0$时，$\ln(1+x)\to0$，所以由夹逼准则可知$\xi\to0$.

          同时，$\frac{\xi}{x}$介于$1$与$\frac{\ln(1+x)}{x}$之间，显然$\lim_{x\to0}\frac{\xi}{x}
              =1$.故
          \[
              \begin{aligned}
                  \lim_{x\to0}\frac{f'(\xi)-f'(0)}{\xi-0}\cdot\frac{\xi}{x}\cdot\frac{x-\ln(1+x)}{x^2} & =f''(0)\cdot\frac{x-\left(x-\frac{x^2}{2}+o(x^2)\right)}{x^2} \\
                                                                                                       & =\frac{f''(0)}{2}.
              \end{aligned}
          \]

    \item \textbf{Solution}. 因为$f(x)$在$x=1$处可导，从而连续.

          在等式$f(1+\sin x)-3f(1-\sin x)=8x+o(x)$两边令$x\to0$得$f(1)=0$，且
          \[
              \begin{aligned}
                  \lim_{x\to0}\frac{f(1+\sin x)-3f(1-\sin x)}{x}= & \lim_{x\to0}\frac{8x+o(x)}{x}=8                                                                                                      \\
                  =                                               & \lim_{x\to0}\frac{f(1+\sin x)-f(1)}{x}-3\lim_{x\to0}\frac{f(1-\sin x)-f(1)}{x}                                                       \\
                  =                                               & \lim_{x\to0}\frac{f(1+\sin x)-f(1)}{\sin x}\cdot\frac{\sin x}{x}-3\lim_{x\to0}\frac{f(1-\sin x)-f(1)}{-\sin x}\cdot\frac{-\sin x}{x} \\
                  =                                               & 4f'(1).
              \end{aligned}
          \]

          所以$f'(1)=2$.

          又$f(x)$的周期为5，所以$f(6)=f(1)=0$，$f'(6)=f'(1)=2$，

          因此曲线$y=f(x)$在点$(6,f(6))$处的切线方程为$y=2(x-6)$即$2x-y-12=0$.

    \item \textbf{Solution}.显然$x_{n}>2$.常数$l=1+\sqrt{2}$满足$l=2+\frac{1}{l}$.下证数列$\{
              x_{n}\}$以$l$为极限.

          \[
              \begin{aligned}
                  0\leq|x_{n+1}-l|= & \left|\left(2+\frac{1}{x_n}\right)-1-\sqrt{2}\right|=\frac{\left|x_n-\left(\sqrt{2}+1\right)\right|}{\left(\sqrt{2}+1\right)x_n} \\
                  \leq              & \frac{1}{2}|x_{n}-l|                                                                                                             \\
                  \leq              & \cdots \leq \frac{1}{2^{n-1}}|x_{2}-l|.
              \end{aligned}
          \]

          由迫敛性知$|x_{n}-l|$收敛到$0$，故数列$\{x_{n}\}$以$l$为极限.
\end{enumerate}

\section{证明题(每小题 5 分，共 10 分)}
\begin{enumerate}
    \setcounter{enumi}{17}
    \item \textbf{Proof}. （1）假设$\forall\,x\in(a,b),f(x)\neq0$，不妨设$f(
              x)>0$，则有
          \[
              \begin{aligned}
                  f'_{+}(a) = & \lim_{x\to a^+}\frac{f(x)-f(a)}{x-a}=\lim_{x\to a^+}\frac{f(x)}{x-a}\geq0, \\
                  f'_{-}(b) = & \lim_{x\to b^-}\frac{f(x)-f(b)}{x-b}=\lim_{x\to b^-}\frac{f(x)}{x-b}\leq0,
              \end{aligned}
          \]

          这与$f'_{+}(a)\cdot f'_{-}(b)>0$矛盾，所以至少存在一点$\xi\in(a,b)$使得$f(\xi
              )=0$.

          （2）由Rolle定理，$\exists\,\eta_{1}\in(a,\xi)$使得$f'(\eta_{1})=0$，$\exists
              \,\eta_{2}\in(\xi,b)$使得$f'(\eta_{2})=0$.

          再由Rolle定理，$\exists\,\eta\in(\eta_{1},\eta_{2})\subset(a,b)$使得$f''(\eta
              )=0$.

          所以至少存在一点$\eta\in(a,b)$使得$f''(\eta)=0$.

    \item \textbf{Proof}. 由Taylor公式，
          \[
              \begin{gathered}
                  f(0)=f(c)+f'(c)(0-c)+\frac{f''(\xi_{1})}{2}(0-c)^2,\quad 0<\xi_1<c<1,\\
                  f(1)=f(c)+f'(c)(1-c)+\frac{f''(\xi_{2})}{2}(1-c)^2,\quad 0<c<\xi_2<1.
              \end{gathered}
          \]

          两式相减，得$f(1)-f(0)=f'(c)+\frac{1}{2}\left[f''(\xi_{2})(1-c)^{2}-f''(\xi
                  _{1})c^{2}\right]$，所以
          \[
              f'(c)=f(1)-f(0)-\frac{1}{2}\left[f''(\xi_{2})(1-c)^{2}-f''(\xi_{1})c^{2}\right
              ].
          \]

          由$|f(x)|\leq a$，$|f''(x)|\leq b$，得$|f'(c)|\leq 2a+\frac{b}{2}\left[(1-c
                  )^{2}+c^{2}\right]\leq 2a+\frac{b}{2}$.
\end{enumerate}
\end{document}



